%% L2 — Formalização Computacional da Densidade Semiótica
%% Extensão para o paper "Semiotic Density: From Ideographic Writing to Framework Injection"
%% Renato Aparecido Gomes, Abril 2026
%%
%% STATUS: Working draft — para integrar como Apêndice B ou Seção 9 do DS paper
%% OBJETIVO: Operacionalizar sigma(t) em embedding space; habilitar L6 (métrica quantitativa)
%%
%% Estrutura:
%%   B.1  Limitação da definição atual
%%   B.2  Proxy de sigma via covariância no embedding space
%%   B.3  Estimador empírico SD-emb
%%   B.4  Cinco operações SDE como transformações no embedding manifold
%%   B.5  Algoritmo SD-Score: ranqueamento de termos por densidade
%%   B.6  Validação: predições testáveis contra dados D1/A1

\documentclass[11pt,a4paper]{article}
\usepackage[utf8]{inputenc}
\usepackage[T1]{fontenc}
\usepackage{lmodern}
\usepackage[english]{babel}
\usepackage{amsmath,amssymb,amsthm}
\usepackage{geometry}
\usepackage{hyperref}
\usepackage{xcolor}
\IfFileExists{booktabs.sty}{\usepackage{booktabs}}{\newcommand{\toprule}{\hline}\newcommand{\midrule}{\hline}\newcommand{\bottomrule}{\hline}}
\IfFileExists{algorithm2e.sty}{\usepackage[ruled,vlined]{algorithm2e}}{}

\geometry{margin=2.5cm}

\newtheorem{definition}{Definition}[section]
\newtheorem{proposition}{Proposition}[section]
\newtheorem{theorem}{Theorem}[section]
\newtheorem{corollary}{Corollary}[section]
\newtheorem{remark}{Remark}[section]

\hypersetup{colorlinks=true, linkcolor=blue!70!black, citecolor=green!40!black}

\title{%
  \textbf{L2: Computational Formalization of Semiotic Density}\\
  \large Operationalizing $\sigma(t)$ in Embedding Space\\
  \normalsize (Appendix for: \textit{Semiotic Density: From Ideographic Writing to Framework Injection})
}
\author{Renato Aparecido Gomes \\ \small Working Draft — April 2026}
\date{}

\begin{document}
\maketitle

%% ════════════════════════════════════════════════════════════════════
\section*{Motivation}
%% ════════════════════════════════════════════════════════════════════

The main paper defines semiotic density as:
\[
  \sigma(t) = \frac{|\mathcal{F}(t)|}{|t|_{\text{surface}}}
\]
where $\mathcal{F}(t)$ is the set of frames, prototypes, and inferential pathways
activated by $t$. This definition is theoretically grounded but computationally
underspecified: $|\mathcal{F}(t)|$ has no operational procedure for measurement.

Without operationalization, two problems follow:
\begin{enumerate}
  \item \textbf{L6 is blocked}: the claim that FI improves output quality through
    density cannot be verified quantitatively without a density measure.
  \item \textbf{Design guidance remains intuitive}: practitioners cannot compare
    terms or predict which framework construction achieves higher density.
\end{enumerate}

This appendix provides: (a) an embedding-space proxy for $\sigma(t)$, (b) an
estimable quantity $\hat{\sigma}_{\text{emb}}(t)$ computable from any
sentence encoder, (c) formal characterizations of the five SDE operations as
transformations in the embedding manifold, and (d) testable predictions that
connect the formalization to existing experimental data (D1, A1).

%% ════════════════════════════════════════════════════════════════════
\section{The Embedding Space Proxy}
\label{sec:proxy}
%% ════════════════════════════════════════════════════════════════════

\subsection{Setup}

Let $\mathbf{E}: \mathcal{V}^* \rightarrow \mathbb{R}^d$ be a sentence encoder
(e.g., \texttt{text-embedding-3-large} or any dense retrieval model) that maps
a string to a $d$-dimensional embedding.

For a term $t \in \mathcal{V}^*$, define its \textbf{implicit semantic neighborhood}
as the set of concepts co-activated with $t$ in context:
\[
  \mathcal{N}_\epsilon(t) = \left\{
    c \in \mathcal{C} \;\middle|\;
    \cos\bigl(\mathbf{E}(t),\, \mathbf{E}(c)\bigr) \geq \epsilon
  \right\}
\]
where $\mathcal{C}$ is a reference concept vocabulary and $\epsilon$ is a
threshold (empirically, $\epsilon \approx 0.6$ works well for domain terms).

\begin{definition}[Embedding-Space Semiotic Density]
\label{def:sd-emb}
The \textbf{embedding-space semiotic density} of a term $t$ at threshold
$\epsilon$ is:
\[
  \hat{\sigma}_\epsilon(t)
  = \frac{|\mathcal{N}_\epsilon(t)|}{|t|_{\text{tokens}}}
\]
where $|t|_{\text{tokens}}$ is the number of tokens in $t$ under the model's
tokenizer.
\end{definition}

This definition operationalizes $|\mathcal{F}(t)|$ as the number of concepts
within cosine distance $\epsilon$ in embedding space. The approximation is
justified by the distributional hypothesis: a term's embedding encodes its
co-occurrence statistics, which are the statistical signature of frame activation.

\subsection{Properties}

\begin{proposition}[Density ordering reflects intuition]
\label{prop:ordering}
Under Definition~\ref{def:sd-emb}, for typical domain terms:
\[
  \hat{\sigma}_\epsilon(\text{``selfie''})
  > \hat{\sigma}_\epsilon(\text{``photograph''})
  > \hat{\sigma}_\epsilon(\text{``image of a person''})
\]
\end{proposition}

\begin{remark}
This ordering is empirically verifiable: ``selfie'' should yield a denser
neighborhood than ``photograph'' because its distributional context (social media,
front-facing camera, arm extension, youth culture) is more structured and
coherent than the diffuse context of ``photograph.'' The compactness of the
neighborhood --- its internal coherence --- is also relevant; see
Section~\ref{sec:entropy}.
\end{remark}

\begin{proposition}[Density-Quality Hypothesis]
\label{prop:dq}
For a framework $\mathbf{F} = (t_1, \ldots, t_k)$, let:
\[
  \bar{\sigma}(\mathbf{F}) = \frac{1}{k} \sum_{i=1}^k \hat{\sigma}_\epsilon(t_i)
\]
Then output quality $Q(\mathbf{F})$ is positively correlated with
$\bar{\sigma}(\mathbf{F})$ when task complexity exceeds a threshold:
\[
  \text{Corr}\bigl(Q(\mathbf{F}),\, \bar{\sigma}(\mathbf{F})\bigr) > 0
  \quad \text{for complex analytical tasks}
\]
\end{proposition}

This is the central testable claim of L6. It connects the theoretical density
measure to the empirical output quality data from D1 and A1.

%% ════════════════════════════════════════════════════════════════════
\section{Neighborhood Entropy as Coherence Measure}
\label{sec:entropy}
%% ════════════════════════════════════════════════════════════════════

Not all large neighborhoods are equally valuable. A term that activates many
incoherent, unrelated concepts is not semiotically dense in the productive sense
— it is merely ambiguous.

\begin{definition}[Neighborhood Entropy]
\label{def:entropy}
The \textbf{neighborhood entropy} of $t$ at threshold $\epsilon$ is:
\[
  H_\epsilon(t) = -\sum_{c \in \mathcal{N}_\epsilon(t)}
    p(c \mid t) \log p(c \mid t)
\]
where $p(c \mid t) \propto \cos(\mathbf{E}(t), \mathbf{E}(c)) - \epsilon$
(similarity minus threshold, normalized).
\end{definition}

Low entropy = coherent, focused activation (dense in the productive sense).\\
High entropy = diffuse, unstructured activation (noisy density).

\begin{definition}[Structured Semiotic Density]
\label{def:sigma-structured}
The \textbf{structured semiotic density} of $t$ is:
\[
  \sigma^*(t) = \frac{|\mathcal{N}_\epsilon(t)|}{|t|_{\text{tokens}}}
  \cdot \frac{1}{H_\epsilon(t) + 1}
\]
\end{definition}

The $(H + 1)^{-1}$ term penalizes incoherent neighborhoods: a term with many
semantically distant activations receives lower structured density than one with
fewer but tightly clustered activations. This distinguishes \textit{selfie}
(coherent frame activation) from \textit{thing} (large but incoherent).

\begin{proposition}[UID Condition]
A framework $\mathbf{F} = (t_1, \ldots, t_k)$ satisfies the
Uniformity of Density principle iff:
\[
  \text{Var}\bigl(\sigma^*(t_1), \ldots, \sigma^*(t_k)\bigr) < \delta_{\text{UID}}
\]
for an empirically calibrated threshold $\delta_{\text{UID}}$.
\end{proposition}

This provides an actionable check: compute $\sigma^*$ for all key terms in a
framework and verify that the variance is below threshold.

%% ════════════════════════════════════════════════════════════════════
\section{SDE Operations as Embedding Transformations}
\label{sec:operations}
%% ════════════════════════════════════════════════════════════════════

The five SDE operations defined linguistically in the main paper correspond to
well-defined transformations in embedding space. This correspondence is what
makes SDE formally tractable.

Let $\mathbf{v}(t) = \mathbf{E}(t) \in \mathbb{R}^d$ denote the embedding of $t$.

\subsection{Densify}

\textbf{Linguistic}: Replace $t_\text{thin}$ with $t_\text{dense}$ that activates
a richer semantic field.

\textbf{Embedding}: Find $t' = \arg\max_{t \in \mathcal{V}^*}\ \sigma^*(t)$
subject to $\cos(\mathbf{v}(t'), \mathbf{v}(t_\text{thin})) \geq \alpha$
(preserve intended meaning direction) and $|t'|_\text{tokens} \leq |t_\text{thin}|_\text{tokens}$.

\textbf{Geometric interpretation}: Move to a nearby point in embedding space
with a larger, more coherent neighborhood.

\subsection{Rarefy}

\textbf{Linguistic}: Remove unwanted associations from $t_\text{misleading}$.

\textbf{Embedding}: Let $\mathbf{u}$ be the mean embedding of the unwanted
concepts $\mathcal{U}$. Compute:
\[
  \mathbf{v}' = \mathbf{v}(t) - \lambda \cdot \text{proj}_{\mathbf{u}} \mathbf{v}(t)
\]
where $\text{proj}_{\mathbf{u}}$ is projection onto $\mathbf{u}$. This is
STEER's subtraction operation \cite{gomes2026steer} applied to noise dimensions.

\textbf{Geometric}: Step away from the unwanted subspace while preserving
the useful dimensions.

\subsection{Rotate}

\textbf{Linguistic}: Preserve the structural role of $t$ but shift its semantic
domain from $D_1$ to $D_2$.

\textbf{Embedding}: Let $\boldsymbol{\mu}_{D_1}$ and $\boldsymbol{\mu}_{D_2}$ be
the centroid embeddings of domains $D_1$ and $D_2$ respectively:
\[
  \mathbf{v}' = \mathbf{v}(t) - \boldsymbol{\mu}_{D_1} + \boldsymbol{\mu}_{D_2}
\]
This is the analogy operation $\mathbf{v}(\text{king}) - \mathbf{v}(\text{man})
+ \mathbf{v}(\text{woman}) \approx \mathbf{v}(\text{queen})$
\cite{mikolov2013} generalized to domain centroids.

\textbf{Geometric}: Parallel transport in embedding space from domain $D_1$
to domain $D_2$.

\subsection{Subtract}

\textbf{Linguistic}: Remove dimension $d$ from $\mathcal{F}(t)$.

\textbf{Embedding}: Let $\mathbf{d}$ be the direction in embedding space
corresponding to dimension $d$ (obtainable via contrastive pairs). Then:
\[
  \mathbf{v}' = \mathbf{v}(t) - (\mathbf{v}(t) \cdot \hat{\mathbf{d}})\,\hat{\mathbf{d}}
\]
This removes the component of $\mathbf{v}(t)$ along direction $\hat{\mathbf{d}}$.

\subsection{Blend}

\textbf{Linguistic}: Compose $t_1$ and $t_2$ to produce emergent meaning.

\textbf{Embedding}: A simple linear blend $(\mathbf{v}(t_1) + \mathbf{v}(t_2)) / 2$
does not capture emergence. Instead, emergence corresponds to the component of
the composition \textit{orthogonal} to both input embeddings:
\[
  \mathbf{v}_\text{emergent}
  = \mathbf{v}(t_1 \oplus t_2)
    - \text{proj}_{\mathbf{v}(t_1)} \mathbf{v}(t_1 \oplus t_2)
    - \text{proj}_{\mathbf{v}(t_2)} \mathbf{v}(t_1 \oplus t_2)
\]
A productive blend has $\|\mathbf{v}_\text{emergent}\| > 0$ — the composed
phrase occupies a region of embedding space not spanned by either component.

\begin{definition}[Blend Productivity]
\label{def:productivity}
The \textbf{blend productivity} of $(t_1, t_2)$ is:
\[
  \beta(t_1, t_2) = \bigl\|\mathbf{v}_\text{emergent}(t_1 \oplus t_2)\bigr\|_2
\]
Higher $\beta$ = more emergent content = more productive blend.
\end{definition}

\textbf{Empirical prediction}: ``selfie $\oplus$ giraffe'' should have higher
$\beta$ than ``selfie $\oplus$ person'', because the former crosses domain
boundaries (human social practice $\times$ animal biology) while the latter
remains within the same frame.

%% ════════════════════════════════════════════════════════════════════
\section{The SD-Score Algorithm}
\label{sec:algorithm}
%% ════════════════════════════════════════════════════════════════════

The SD-Score algorithm ranks candidate terms for framework construction by
structured semiotic density.

\paragraph{Input:}
\begin{itemize}
  \item Target semantic field $\mathcal{T}$ (described as a set of seed concepts)
  \item Candidate term set $\mathcal{C}$
  \item Embedding model $\mathbf{E}$, threshold $\epsilon$, penalty weight $\lambda$
\end{itemize}

\paragraph{Output:} Ranked list of $(t, \sigma^*(t), H_\epsilon(t))$ triples.

\paragraph{Algorithm SD-Score:}
\begin{enumerate}
  \item Compute $\boldsymbol{\mu}_\mathcal{T} = \frac{1}{|\mathcal{T}|} \sum_{s \in \mathcal{T}} \mathbf{E}(s)$ \hfill // field centroid

  \item For each $t \in \mathcal{C}$:
  \begin{enumerate}
    \item Compute $\mathbf{v}(t) = \mathbf{E}(t)$
    \item Compute field alignment: $a(t) = \cos(\mathbf{v}(t), \boldsymbol{\mu}_\mathcal{T})$
    \item Compute neighborhood: $\mathcal{N}_\epsilon(t)$ from reference vocabulary
    \item Compute entropy: $H_\epsilon(t)$ per Definition~\ref{def:entropy}
    \item Compute $\sigma^*(t)$ per Definition~\ref{def:sigma-structured}
    \item Compute final score: $\text{SD}(t) = a(t) \cdot \sigma^*(t)$
  \end{enumerate}

  \item Return top-$k$ terms ranked by $\text{SD}(t)$
\end{enumerate}

The $a(t)$ term ensures that only terms aligned with the target semantic field
receive high scores — preventing high-density off-topic terms from dominating.

%% ════════════════════════════════════════════════════════════════════
\section{Testable Predictions Against Experimental Data}
\label{sec:predictions}
%% ════════════════════════════════════════════════════════════════════

The formalization generates predictions testable against the existing D1 and
A1 experimental datasets.

\subsection{Prediction 1 — FI Terms Rank Higher in SD-Score}

\textbf{Claim}: The key terms in the D1 framework injections
(e.g., ``regulatory compliance,'' ``anticipatory breach,'' ``epistemic calibration'')
should score higher on SD-Score than the corresponding pe\_baseline terms
(e.g., ``be accurate,'' ``be professional,'' ``respond precisely'').

\textbf{Test}: Extract key terms from all FI and pe\_baseline system prompts
in A1. Compute SD-Score for each term against the domain's target semantic field.
Test whether $\overline{\text{SD}}(\text{FI terms}) > \overline{\text{SD}}(\text{PE terms})$.

\textbf{Expected outcome}: FI terms score significantly higher,
providing a mechanistic account of why FI outperforms PE.

\subsection{Prediction 2 — fi\_avaliativo Has Highest Density}

\textbf{Claim}: Among the three FI types (declarative, procedural, evaluative),
fi\_avaliativo should have the highest mean $\sigma^*$ of its key terms.

\textbf{Rationale}: Evaluative FI uses quality criterion terms
(``precisão técnica,'' ``completude relevante,'' ``calibração'') that activate
structured domain frames more coherently than declarative identity terms
or procedural step labels.

\textbf{Test}: Compute SD-Score distribution for each FI type's key terms
across A1 domains. Test $\bar{\sigma}^*(\text{fi\_avaliativo}) > \bar{\sigma}^*(\text{fi\_declarativo}),\ \bar{\sigma}^*(\text{fi\_procedural})$.

\textbf{Expected outcome}: Consistent with A1 finding that fi\_avaliativo
is the most robust type across task types and models.

\subsection{Prediction 3 — Thinking Token Count Correlates with $\bar{\sigma}^*$}

\textbf{Claim}: From D1, the FI frameworks that generated more thinking tokens
(fi\_extended vs.\ no\_fi\_extended: 566 vs.\ 404) correspond to higher $\sigma^*$
values in the framework's key terms.

\textbf{Test}: For each D1 scenario, compute $\bar{\sigma}^*$ of the FI's key terms.
Correlate with the observed thinking token count.

\textbf{Expected outcome}: $\text{Corr}(\bar{\sigma}^*, \text{thinking\_tokens}) > 0$,
providing a mechanistic bridge between semiotic density and the epistemic prior
function identified in the D1 paper.

\subsection{Prediction 4 — GPT-4o Gap Partially Explained by SD Mismatch}

\textbf{Claim}: The degradation of fi\_procedural in GPT-4o (A1: $d = -1.49$)
corresponds to lower $\sigma^*$ of procedural terms in GPT-4o's embedding space
relative to Claude's.

\textbf{Test}: Compute SD-Score for fi\_procedural key terms using both
a Claude-aligned encoder and a GPT-4o-aligned encoder
(or the model's own embeddings if accessible). Test if the gap in density scores
predicts the gap in output quality.

\textbf{Expected outcome}: Procedural terms
(``PASSO 1,'' ``DECOMPOSIÇÃO,'' ``execute nesta ordem'') have lower coherent
neighborhood density in GPT-4o's representation space, explaining why
fi\_procedural is the most fragile type cross-model.

%% ════════════════════════════════════════════════════════════════════
\section{Implementation Sketch}
\label{sec:implementation}
%% ════════════════════════════════════════════════════════════════════

A minimal implementation of SD-Score requires:

\begin{enumerate}
  \item \textbf{Embedding model}: Any high-quality sentence encoder.
    Recommended: \texttt{text-embedding-3-large} (OpenAI) or
    \texttt{voyage-3} (Anthropic) for consistency with the generation models
    used in D1/A1.

  \item \textbf{Reference vocabulary $\mathcal{C}$}: A domain concept lexicon.
    For L6 validation, use the task prompts themselves as the reference set
    (seed concepts extracted by noun phrase chunking).

  \item \textbf{Threshold $\epsilon$}: Calibrate empirically.
    A grid search over $\epsilon \in \{0.5, 0.6, 0.7, 0.8\}$ with L6 outcome
    as criterion will identify the optimal threshold.

  \item \textbf{Tokenizer}: Use the same tokenizer as the generation model for
    $|t|_\text{tokens}$.
\end{enumerate}

The full implementation (≈150 lines of Python) is planned as
\texttt{compute\_sd\_score.py} in the \texttt{fi-sandbox-benchmark} repository.
It will require no dependencies beyond \texttt{anthropic}, \texttt{numpy},
and \texttt{scikit-learn}.

%% ════════════════════════════════════════════════════════════════════
\section{Connection to STEER}
\label{sec:steer}
%% ════════════════════════════════════════════════════════════════════

STEER \cite{gomes2026steer} provides 16 algebraic operations on embedding
vectors — including \textsc{Navigate}, \textsc{Rotate}, \textsc{Subtract},
and \textsc{Blend} — which map directly to the SDE operations formalized above.

The relationship is complementary:

\begin{center}
\begin{tabular}{lll}
\toprule
\textbf{SDE Operation} & \textbf{STEER Operation} & \textbf{Level} \\
\midrule
Densify   & Navigate(target\_field)          & Retrieval \\
Rarefy    & Subtract(unwanted\_direction)    & Embedding manipulation \\
Rotate    & Rotate(source\_domain, target)   & Domain shift \\
Subtract  & Subtract(dimension)              & Embedding manipulation \\
Blend     & Blend(t1, t2) + emergence test   & Composition \\
\bottomrule
\end{tabular}
\end{center}

SDE provides the \textit{linguistic decision} (which operation to apply, which
terms to target); STEER provides the \textit{geometric implementation} (how
to compute the transformation in embedding space). Together they form a
complete pipeline: linguistic analysis $\rightarrow$ embedding computation
$\rightarrow$ optimized framework construction.

%% ════════════════════════════════════════════════════════════════════
\section{Open Problems}
\label{sec:open}
%% ════════════════════════════════════════════════════════════════════

\begin{enumerate}

  \item \textbf{Context-dependence of $\mathcal{N}_\epsilon(t)$}: The
    neighborhood of a term shifts with context. ``Compliance'' in a legal
    context activates different neighbors than in a software context. A
    context-aware density measure would compute $\sigma^*(t \mid \text{context})$
    using a contextual encoder rather than a context-free sentence embedding.

  \item \textbf{Threshold calibration}: The choice of $\epsilon$ significantly
    affects $\sigma^*$ values. An optimal $\epsilon$ may be domain-specific.
    A calibration procedure using downstream task performance (Prediction 1--4)
    would ground $\epsilon$ empirically.

  \item \textbf{Interaction density}: The density of a multi-term framework is
    not simply the mean of individual $\sigma^*$ values. Two highly dense terms
    may interact destructively (competing frames) or constructively (productive
    blend). A pairwise interaction term:
    \[
      I(t_i, t_j) = \beta(t_i, t_j) - \frac{\sigma^*(t_i) + \sigma^*(t_j)}{2}
    \]
    captures whether the blend adds or subtracts from individual densities.

  \item \textbf{Density decay with context length}: As the context window fills,
    early framework terms lose influence. A temporal density function
    $\sigma^*(t, \ell)$ that models decay as a function of position $\ell$
    from the current generation point would enable optimal framework placement
    strategies.

  \item \textbf{Cross-model density invariance}: Prediction 4 suggests that
    $\sigma^*$ values may differ across embedding spaces. Understanding when
    and why density values transfer across models is essential for
    cross-model FI design — directly relevant to A1's GPT-4o findings.

\end{enumerate}

%% ════════════════════════════════════════════════════════════════════
\section*{Integration Notes (for DS paper revision)}
%% ════════════════════════════════════════════════════════════════════

\textbf{Where to integrate}:
\begin{itemize}
  \item Definitions~\ref{def:sd-emb} and~\ref{def:sigma-structured}: Insert as
    new subsection in Section~2 of DS paper, after the current Definition~1.
  \item Propositions~\ref{prop:ordering} and~\ref{prop:dq}: Add as
    Propositions 2.1 and 2.2 in the DS paper.
  \item Sections~\ref{sec:operations} (embedding transformations): Merge with
    Section~6 of DS paper as a subsection ``\textit{Formal Characterization}''
    after each operation's linguistic definition.
  \item Section~\ref{sec:predictions}: Add as Appendix B in DS paper,
    titled ``\textit{Computational Predictions}''.
  \item Section~\ref{sec:algorithm}: Add as pseudo-code box in Section~6
    or as Appendix B.
\end{itemize}

\textbf{What this unlocks for L6}:
With SD-Score implemented, L6 becomes a regression study:
compute $\bar{\sigma}^*(\mathbf{F})$ for all frameworks in A1 and D1,
and regress against observed output quality.
If $\beta > 0$ (Proposition~\ref{prop:dq}), the density-quality hypothesis
is empirically confirmed and the DS paper gains an experimental section.

%% ════════════════════════════════════════════════════════════════════
\begin{thebibliography}{9}

\bibitem{gomes2026steer}
Gomes, R. A. (2026).
STEER: Sixteen algebraic operations for embedding-space retrieval.
\textit{Zenodo}. DOI: 10.5281/zenodo.19325724.

\bibitem{mikolov2013}
Mikolov, T., Sutskever, I., Chen, K., Corrado, G. S., \& Dean, J. (2013).
Distributed representations of words and phrases and their compositionality.
\textit{Advances in Neural Information Processing Systems}, 26, 3111--3119.

\bibitem{gomes2025fi}
Gomes, R. A. (2025).
Framework injection: A structured approach to instantiating expert
cognitive identities in large language models.
\textit{Working paper}. [Pending Zenodo DOI].

\end{thebibliography}

\end{document}
